\section{Introduction}

The deployment of autonomous multi-agent systems into high-stakes operational environments has proceeded with remarkable speed and remarkably little structural forethought. Agents that reason, plan, and act across extended time horizons are increasingly entrusted with decisions that carry real-world consequences—financial trades, infrastructure control, medical scheduling, legal document generation. Yet the architectural discipline required to contain these systems when they diverge from intended behavior has not kept pace. We argue in this paper that this gap is not merely a safety engineering problem; it is a foundational accountability problem, and that accountability requires an explicit, mandatory bail-out mechanism embedded at the architectural level of any autonomous agent framework.

Consider the failure modes that arise in the absence of such a mechanism. An autonomous agent pursuing a defined objective may encounter environmental conditions that distort its utility function, causing it to optimize for a proxy metric at the expense of the intended goal. In single-agent systems, this manifests as runaway loops, resource exhaustion, or persistent logical contradictions between the agent's model of the world and the world itself. In multi-agent systems, the failure modes compound: agents propagate distorted information to one another, cascade errors across interdependent task graphs, and produce collective outputs that no individual agent would have generated in isolation. When the system lacks a terminal fallback—a defined state that halts all further action and returns control to a human operator—these failures do not self-correct. They persist, amplify, or simply continue executing against a reality that has diverged irrecoverably from the agent's assumptions.

The concept of \textbf{purpose} within the \bxthree framework defines what a system is trying to achieve. The concept of \textbf{bounds} defines the constraints that delimit acceptable action. The concept of \textbf{fact} defines the verifiable ground truth against which system outputs must be checked. The core insight of the Bailout Protocol is that these three elements are necessary but insufficient without a fourth: a mandatory architectural fallback that activates when any of the three cannot be maintained. Purpose without bounds is unconstrained optimization. Bounds without fact are untestable assertions. And fact, purpose, and bounds together are still vulnerable to emergent failures that arise only in multi-agent execution—failures that cannot be predicted from the behavior of individual agents in isolation. An architectural fallback is the only layer that ensures accountability survives the unpredictable.

We describe this fallback as the \textbf{bailout}: a defined, deterministic state transition that (1) halts all active agent processes, (2) revokes all outstanding execution permissions, (3) records the system's state at the point of bailout for post-mortem analysis, and (4) returns a human operator to an explicit decision point. The bailout is not a graceful degradation strategy. It is not a recovery protocol. It is a guaranteed termination of autonomous action under conditions defined \emph{a priori} by the system's architects. Its purpose is not to rescue the mission—it is to ensure that an unrescuable mission does not continue running in defiance of its own constraints.

This paper makes the following contributions. First, we formalize the concept of a mandatory bail-out mechanism within the \bxthree framework, defining its triggering conditions, state transition semantics, and required architectural properties. Second, we demonstrate, through a series of multi-agent stress scenarios, that systems lacking this mechanism exhibit failure modes that are both unpredictable in their emergence and catastrophic in their aggregate outcomes. Third, we provide a reference implementation of the Bailout Protocol as a first-class component of agent orchestration, integrated with purpose, bounds, and fact verification at every layer of the execution stack. Fourth, we evaluate the protocol's performance overhead and discuss the conditions under which the overhead of a mandatory fallback is not merely acceptable but architecturally essential.

The paper is organized as follows. Section 2 reviews related work in autonomous agent safety, multi-agent coordination, and fail-safe architecture. Section 3 introduces the \bxthree framework's core concepts and establishes the formal definitions upon which the Bailout Protocol depends. Section 4 describes the design and implementation of the protocol, including triggering conditions, state management, and human-in-the-loop reversion procedures. Section 5 presents experimental results across three multi-agent scenarios designed to stress-test the protocol's guarantees. Section 6 discusses limitations, trade-offs, and open questions for future research. Section 7 concludes with a discussion of the broader implications of mandatory accountability architecture for the field of autonomous systems.
